\chapter{Base de données Supabase}

La base est une instance \textbf{PostgreSQL} gérée par Supabase. Elle comporte
\textbf{21 tables} dans le schéma \texttt{public}, protégées par RLS. Cette
section décrit chaque table, ses colonnes principales et son rôle. Les
énumérations utilisées sont~:
\begin{itemize}
  \item \texttt{app\_role} : \texttt{student, parent, admin, pedago} ;
  \item \texttt{school\_level} : \texttt{5eme\_primaire, 1ere\_cem, 2eme\_cem,
        3eme\_cem, 4eme\_cem, premiere, seconde, terminale} ;
  \item \texttt{link\_status} : \texttt{pending, active, rejected}.
\end{itemize}

\begin{notebox}
Les relations entre tables sont \emph{logiques} (par \texttt{user\_id},
\texttt{chapter\_id}, \texttt{lesson\_id}\dots) et non matérialisées par des clés
étrangères strictes vers \texttt{auth.users} (recommandation Supabase). Le MCD du
chapitre~\ref{chap:mcd} les illustre.
\end{notebox}

% ---------- Macro pour décrire une table ----------
\newcommand{\dbtable}[2]{%
\subsection{\texttt{#1}}%
\emph{#2}%
}

\section{Utilisateurs et rôles}

\dbtable{profiles}{Profil applicatif de chaque utilisateur (extension de auth.users).}
\begin{longtable}{p{3.6cm}p{2.6cm}p{8cm}}
\toprule \textbf{Colonne} & \textbf{Type} & \textbf{Description} \\ \midrule \endhead
\texttt{id} & uuid & Identifiant (= id auth). \\
\texttt{email} & text & Adresse e-mail. \\
\texttt{email\_verified} & boolean & E-mail confirmé. \\
\texttt{first\_name / last\_name} & text & Nom et prénom. \\
\texttt{phone} & text & Téléphone. \\
\texttt{avatar\_url} & text & Photo de profil (bucket \texttt{avatars}). \\
\texttt{school\_level} & enum & Niveau scolaire. \\
\texttt{filiere} & text & Filière (secondaire). \\
\texttt{linking\_code} & text & Code de liaison parent--enfant. \\
\texttt{is\_active} & boolean & Compte actif. \\
\texttt{date\_of\_birth} & date & Date de naissance. \\
\texttt{wilaya / ville / ecole} & text & Localisation et établissement. \\
\bottomrule
\end{longtable}

\dbtable{user\_roles}{Rôles applicatifs (séparés du profil pour la sécurité).}
\begin{longtable}{p{3.6cm}p{2.6cm}p{8cm}}
\toprule \textbf{Colonne} & \textbf{Type} & \textbf{Description} \\ \midrule \endhead
\texttt{id} & uuid & Identifiant. \\
\texttt{user\_id} & uuid & Utilisateur concerné. \\
\texttt{role} & app\_role & Rôle (student/parent/admin/pedago). \\
\texttt{created\_at} & timestamptz & Date d'attribution. \\
\bottomrule
\end{longtable}

\dbtable{parent\_child\_links}{Liaison entre un parent et un enfant.}
\begin{longtable}{p{3.6cm}p{2.6cm}p{8cm}}
\toprule \textbf{Colonne} & \textbf{Type} & \textbf{Description} \\ \midrule \endhead
\texttt{parent\_id} & uuid & Parent. \\
\texttt{child\_id} & uuid & Enfant. \\
\texttt{status} & link\_status & État (pending/active/rejected). \\
\bottomrule
\end{longtable}

\dbtable{activity\_logs}{Journal d'audit des actions (via fonction \texttt{log\_activity}).}
\begin{longtable}{p{3.6cm}p{2.6cm}p{8cm}}
\toprule \textbf{Colonne} & \textbf{Type} & \textbf{Description} \\ \midrule \endhead
\texttt{user\_id} & uuid & Auteur de l'action. \\
\texttt{action} & text & Code de l'action. \\
\texttt{details} & jsonb & Détails contextuels. \\
\texttt{ip\_address} & text & Adresse IP (optionnelle). \\
\bottomrule
\end{longtable}

\section{Contenu pédagogique}

\dbtable{filieres}{Filières disponibles par niveau scolaire.}
\begin{longtable}{p{3.6cm}p{2.6cm}p{8cm}}
\toprule \textbf{Colonne} & \textbf{Type} & \textbf{Description} \\ \midrule \endhead
\texttt{school\_level} & enum & Niveau associé. \\
\texttt{code} & text & Code court (ex.\ \texttt{sciences}). \\
\texttt{name / name\_ar} & text & Libellé FR / AR. \\
\bottomrule
\end{longtable}

\dbtable{chapters}{Chapitres de cours, organisés par niveau, filière et matière.}
\begin{longtable}{p{3.6cm}p{2.6cm}p{8cm}}
\toprule \textbf{Colonne} & \textbf{Type} & \textbf{Description} \\ \midrule \endhead
\texttt{school\_level} & enum & Niveau scolaire. \\
\texttt{filiere\_id} & uuid & Filière (nullable). \\
\texttt{title / title\_ar} & text & Titre FR / AR. \\
\texttt{description} & text & Description. \\
\texttt{subject} & text & Matière (\texttt{math}). \\
\texttt{order\_index} & integer & Ordre d'affichage. \\
\bottomrule
\end{longtable}

\dbtable{lessons}{Leçons appartenant à un chapitre.}
\begin{longtable}{p{3.6cm}p{2.6cm}p{8cm}}
\toprule \textbf{Colonne} & \textbf{Type} & \textbf{Description} \\ \midrule \endhead
\texttt{chapter\_id} & uuid & Chapitre parent. \\
\texttt{title / title\_ar} & text & Titre FR / AR. \\
\texttt{content} & text & Contenu (Markdown/HTML). \\
\texttt{video\_url} & text & Vidéo associée. \\
\texttt{order\_index} & integer & Ordre. \\
\bottomrule
\end{longtable}

\dbtable{chapter\_exercises}{Exercices d'entraînement (réponse libre).}
\begin{longtable}{p{3.6cm}p{2.6cm}p{8cm}}
\toprule \textbf{Colonne} & \textbf{Type} & \textbf{Description} \\ \midrule \endhead
\texttt{chapter\_id / lesson\_id} & uuid & Rattachement. \\
\texttt{title} & text & Titre de l'exercice. \\
\texttt{statement} & text & Énoncé. \\
\texttt{expected\_answer} & text & Réponse attendue. \\
\texttt{accepted\_answers} & jsonb & Variantes acceptées. \\
\texttt{solution} & text & Solution détaillée. \\
\texttt{hint} & text & Indice. \\
\texttt{difficulty} & integer & Difficulté (1--5). \\
\bottomrule
\end{longtable}

\dbtable{chapter\_quizzes}{Questions de quiz (choix multiples).}
\begin{longtable}{p{3.6cm}p{2.6cm}p{8cm}}
\toprule \textbf{Colonne} & \textbf{Type} & \textbf{Description} \\ \midrule \endhead
\texttt{chapter\_id / lesson\_id} & uuid & Rattachement. \\
\texttt{question} & text & Énoncé de la question. \\
\texttt{options} & jsonb & Liste des choix. \\
\texttt{correct\_answer} & text & Bonne réponse. \\
\texttt{explanation} & text & Explication. \\
\texttt{hint} & text & Indice. \\
\texttt{difficulty} & integer & Difficulté (1--5). \\
\bottomrule
\end{longtable}

\dbtable{exams}{Examens / sujets par trimestre.}
\begin{longtable}{p{3.6cm}p{2.6cm}p{8cm}}
\toprule \textbf{Colonne} & \textbf{Type} & \textbf{Description} \\ \midrule \endhead
\texttt{school\_level} & enum & Niveau. \\
\texttt{trimester} & integer & Trimestre. \\
\texttt{title / title\_ar} & text & Titre. \\
\texttt{content} & jsonb & Contenu structuré. \\
\texttt{duration\_minutes} & integer & Durée. \\
\texttt{subject} & text & Matière. \\
\bottomrule
\end{longtable}

\section{Suivi de l'apprentissage}

\dbtable{student\_scores}{Cœur du moteur adaptatif : niveau et statistiques par élève/leçon/chapitre.}
\begin{longtable}{p{4.0cm}p{2.4cm}p{8cm}}
\toprule \textbf{Colonne} & \textbf{Type} & \textbf{Description} \\ \midrule \endhead
\texttt{user\_id} & uuid & Élève. \\
\texttt{chapter\_id / lesson\_id} & uuid & Portée (leçon, chapitre, ou globale si NULL). \\
\texttt{current\_level} & integer & Niveau courant (5--100). \\
\texttt{reading/quiz/exercise\_time\_seconds} & integer & Temps passé par type. \\
\texttt{total\_answers} & integer & Nombre total de réponses. \\
\texttt{correct\_answers} & integer & Réponses correctes. \\
\texttt{accuracy\_rate} & numeric & Taux de réussite (\%). \\
\texttt{streak} & integer & Série de bonnes réponses. \\
\texttt{assessment\_data} & jsonb & Drapeaux d'hésitation, données du test. \\
\texttt{periodic\_advice} & jsonb & Conseils périodiques. \\
\texttt{advice\_seen} & boolean & Conseil consulté. \\
\texttt{report\_first\_shown\_at} & timestamptz & Premier affichage du rapport. \\
\bottomrule
\end{longtable}

\dbtable{ai\_generated\_content}{Contenu (quiz/exercices/révision) généré par IA et calibré par niveau.}
\begin{longtable}{p{4.0cm}p{2.4cm}p{8cm}}
\toprule \textbf{Colonne} & \textbf{Type} & \textbf{Description} \\ \midrule \endhead
\texttt{user\_id} & uuid & Élève destinataire. \\
\texttt{chapter\_id / lesson\_id} & uuid & Rattachement. \\
\texttt{content\_type} & text & \texttt{quiz} / \texttt{exercise} / \texttt{revision}. \\
\texttt{content} & jsonb & Items générés. \\
\texttt{difficulty\_level} & integer & Niveau de calibration (composite). \\
\bottomrule
\end{longtable}

\dbtable{ai\_lesson\_comments}{Commentaires IA personnalisés après une session d'étude.}
\begin{longtable}{p{4.0cm}p{2.4cm}p{8cm}}
\toprule \textbf{Colonne} & \textbf{Type} & \textbf{Description} \\ \midrule \endhead
\texttt{user\_id} & uuid & Élève. \\
\texttt{lesson\_id / chapter\_id} & uuid & Leçon/chapitre concernés. \\
\texttt{level\_before / level\_after / level\_delta} & integer & Évolution du niveau. \\
\texttt{weak\_concepts / strong\_concepts} & jsonb & Lacunes et acquis détectés. \\
\texttt{message} & text & Commentaire (arabe, Markdown). \\
\texttt{link\_url} & text & Lien vers la leçon. \\
\bottomrule
\end{longtable}

\dbtable{parent\_reports}{Rapports de progression destinés aux parents.}
\begin{longtable}{p{4.0cm}p{2.4cm}p{8cm}}
\toprule \textbf{Colonne} & \textbf{Type} & \textbf{Description} \\ \midrule \endhead
\texttt{parent\_id / child\_id} & uuid & Parent et enfant. \\
\texttt{period\_start / period\_end} & timestamptz & Période couverte. \\
\texttt{report\_data} & jsonb & Données brutes. \\
\texttt{global\_success\_rate} & numeric & Taux de réussite global. \\
\texttt{global\_level} & integer & Niveau global. \\
\texttt{strong\_chapters / weak\_chapters} & jsonb & Forces et faiblesses. \\
\texttt{ai\_recommendations / summary} & text & Recommandations IA. \\
\texttt{report\_type} & text & Type de rapport. \\
\bottomrule
\end{longtable}

\section{Assistant IA (chat)}

\dbtable{chat\_conversations}{Historique des conversations avec l'assistant IA.}
\begin{longtable}{p{3.6cm}p{2.6cm}p{8cm}}
\toprule \textbf{Colonne} & \textbf{Type} & \textbf{Description} \\ \midrule \endhead
\texttt{user\_id} & uuid & Propriétaire. \\
\texttt{chapter\_id} & uuid & Contexte (nullable). \\
\texttt{title} & text & Titre de la conversation. \\
\texttt{messages} & jsonb & Fil de messages. \\
\bottomrule
\end{longtable}

\dbtable{chat\_usage}{Quotas d'utilisation quotidienne du chat (messages/images).}
\begin{longtable}{p{3.6cm}p{2.6cm}p{8cm}}
\toprule \textbf{Colonne} & \textbf{Type} & \textbf{Description} \\ \midrule \endhead
\texttt{user\_id} & uuid & Utilisateur. \\
\texttt{usage\_date} & date & Jour. \\
\texttt{message\_count} & integer & Nombre de messages. \\
\texttt{image\_count} & integer & Nombre d'images. \\
\bottomrule
\end{longtable}

\section{Abonnements et paiements}

\dbtable{subscription\_config}{Configuration des offres (prix individuel/famille).}
\begin{longtable}{p{3.6cm}p{2.6cm}p{8cm}}
\toprule \textbf{Colonne} & \textbf{Type} & \textbf{Description} \\ \midrule \endhead
\texttt{plan\_type} & text & Type d'offre. \\
\texttt{price\_single / price\_family} & integer & Tarifs. \\
\texttt{label} & text & Libellé. \\
\texttt{is\_active} & boolean & Offre active. \\
\bottomrule
\end{longtable}

\dbtable{subscription\_periods}{Périodes de validité des abonnements.}
\begin{longtable}{p{3.6cm}p{2.6cm}p{8cm}}
\toprule \textbf{Colonne} & \textbf{Type} & \textbf{Description} \\ \midrule \endhead
\texttt{label} & text & Libellé. \\
\texttt{start\_date / end\_date} & date & Bornes. \\
\texttt{is\_active} & boolean & Période active. \\
\bottomrule
\end{longtable}

\dbtable{student\_subscriptions}{Abonnement actif d'un élève (décompte de jours).}
\begin{longtable}{p{3.6cm}p{2.6cm}p{8cm}}
\toprule \textbf{Colonne} & \textbf{Type} & \textbf{Description} \\ \midrule \endhead
\texttt{user\_id} & uuid & Élève. \\
\texttt{activation\_code\_id} & uuid & Code utilisé. \\
\texttt{plan\_type} & text & Type d'offre. \\
\texttt{total\_days / days\_used} & integer/numeric & Crédit et consommation. \\
\texttt{is\_paused / paused\_at} & boolean/ts & Mise en pause. \\
\texttt{started\_at / last\_tick\_at} & timestamptz & Démarrage et dernier décompte. \\
\bottomrule
\end{longtable}

\dbtable{activation\_codes}{Codes prépayés d'activation d'abonnement.}
\begin{longtable}{p{3.6cm}p{2.6cm}p{8cm}}
\toprule \textbf{Colonne} & \textbf{Type} & \textbf{Description} \\ \midrule \endhead
\texttt{code} & text & Code unique. \\
\texttt{plan\_type} & text & Offre liée. \\
\texttt{is\_family} & boolean & Offre famille. \\
\texttt{status} & text & État (free/used\dots). \\
\texttt{created\_by / used\_by} & uuid & Créateur / utilisateur. \\
\texttt{payment\_id} & uuid & Paiement associé. \\
\bottomrule
\end{longtable}

\dbtable{payments}{Historique des paiements.}
\begin{longtable}{p{3.6cm}p{2.6cm}p{8cm}}
\toprule \textbf{Colonne} & \textbf{Type} & \textbf{Description} \\ \midrule \endhead
\texttt{user\_id} & uuid & Payeur. \\
\texttt{plan\_type / plan\_label} & text & Offre. \\
\texttt{amount} & integer & Montant. \\
\texttt{is\_family / children\_count} & bool/int & Offre famille. \\
\texttt{period\_id} & uuid & Période. \\
\texttt{status} & text & État du paiement. \\
\bottomrule
\end{longtable}

\section{Fonctions de base de données}
\begin{longtable}{p{4.8cm}p{9.6cm}}
\toprule \textbf{Fonction} & \textbf{Rôle} \\ \midrule \endhead
\texttt{has\_role(uid, role)} & Vérifie un rôle (SECURITY DEFINER, anti-récursion RLS). \\
\texttt{is\_parent\_of(parent, child)} & Vérifie une liaison parent--enfant active. \\
\texttt{user\_has\_any\_role(uid)} & Indique si l'utilisateur a au moins un rôle. \\
\texttt{handle\_new\_user()} & Trigger : crée profil + rôle à l'inscription. \\
\texttt{handle\_updated\_at()} & Trigger : met à jour \texttt{updated\_at}. \\
\texttt{log\_activity(uid, action, details)} & Insère une ligne d'audit. \\
\texttt{generate\_activation\_code()} & Génère un code d'activation aléatoire. \\
\texttt{check\_exercise\_answer(id, rep)} & Valide une réponse d'exercice (comparaison normalisée exacte). \\
\texttt{check\_quiz\_answer(id, rep)} & Valide une réponse de quiz. \\
\texttt{get\_student\_exercises / get\_student\_quizzes} & Renvoie les exercices/quiz d'un chapitre/leçon. \\
\bottomrule
\end{longtable}

\begin{infobox}[\faCheckCircle\ Validation exacte des réponses]
La fonction \texttt{check\_exercise\_answer} normalise la saisie (minuscules,
suppression des espaces) puis effectue une \textbf{comparaison exacte} avec la
réponse attendue et les variantes acceptées. Cela évite qu'une saisie partielle
(ex.\ «~1~») soit acceptée pour une réponse «~12~».
\end{infobox}
